\documentclass[11pt,a4paper]{article}
\usepackage[margin=2.3cm]{geometry}
\usepackage{graphicx}
\usepackage{amsmath}
\usepackage[hidelinks]{hyperref}
\usepackage{microtype}
\title{The $F_{ST}$--LD dissociation: the signature that separates empirical
ancestry sorting from a neutral null}
\author{}
\date{Module E working note --- \today}

\begin{document}
\maketitle
\vspace{-2em}

\begin{abstract}\noindent
Across 20 replicate \textit{F.\ polyctena}\,$\times$\,\textit{F.\ aquilonia} hybrid
populations, the empirical data and a recombination-matched, haplodiploid
\emph{neutral null} (real phased founders, no selection, calibrated so the neutral
LD background matches) differ in a way no neutral demography can produce: at
ancestry-informative (high-DiagnosticIndex) loci the empirical data show
\textbf{high differentiation but low linkage disequilibrium}, while the null shows
exactly the reverse --- \textbf{low differentiation but high LD}. This single
dissociation, visible both in aggregate and SNP-by-SNP (Fig.~\ref{fig:diss}), is the
clean signature of divergent, locus-specific selection on ancestry.
\end{abstract}

\section*{The observation}
We measure two things per SNP and bin them by DiagnosticIndex (DI, how strongly a
marker distinguishes the parental species): among-population $F_{ST}$, and LD between
a SNP and its \emph{directly adjacent} SNP(s) (mean $r^2$; i.e.\ $ld_w$ at a one-SNP
window). The null is the 20 independent simulated demes, each downsampled to the
matching empirical population size.

\begin{figure}[htbp]\centering
\includegraphics[width=\textwidth]{../Figures/moduleE_emp_vs_sim_summary.pdf}
\caption{\textbf{Empirical vs neutral null.} \textbf{(A)} Among-population $F_{ST}$
rises with DI in the data but is flat in the null (drift is DI-blind); the null even
sits slightly \emph{above} the empirical at neutral loci (it mildly over-differentiates
the background), then the empirical pulls far ahead at diagnostic loci. \textbf{(B)}
Adjacent-SNP LD does the opposite: it \emph{rises} with DI in the null (admixture LD)
while the empirical stays low. \textbf{(C--E)} Per-SNP distributions (violin $+$ quartile
box) split by DI class (columns: neutral $\le-90$, mid, diagnostic $>-25$) for the three
per-SNP metrics (rows), over all high-recombination SNPs. At \emph{neutral} loci the
empirical and null distributions coincide on all three --- the calibrated control. The
dissociation opens only at mid/high DI: empirical $F_{ST}$ jumps to $\sim$0.30 while the
null stays $\sim$0.08 (C); adjacent-SNP LD stays low in the empirical but rises in the
null, strongest at diagnostic loci (median 0.06 vs 0.27; D); and empirical sorting
spreads widely in \emph{mixed directions} while the null sits at zero (E).}
\label{fig:diss}
\end{figure}

\section*{Why this is diagnostic of selection, not demography}
\begin{itemize}\itemsep2pt
\item \textbf{The two observations are one phenomenon.} At a diagnostic locus,
populations sort ancestry \emph{divergently and locus-specifically} --- different
populations fix different loci. That builds high per-locus $F_{ST}$ and, because the
\emph{set} of differentiated populations shuffles from one SNP to the next, the
among-population covariance between neighbouring SNPs stays near zero: high $F_{ST}$
inflates the $r^2$ denominator while the numerator does not, so the pooled LD
collapses (within-population LD $\approx0.23$ falls to $\approx0.03$ when pooled).
\item \textbf{The null cannot do both at once.} With no selection, ancestry stays
admixed and genome-wide-correlated, giving strong \emph{admixture} LD at diagnostic
loci (panel B rises) --- but drift cannot differentiate ancestry-informative from
uninformative markers, so $F_{ST}$ is flat across DI (panel A). Every neutral model
(one panmictic pool, 20 independent demes, or a founding tree) is DI-blind and gives
this flat-$F_{ST}$/high-LD profile.
\item \textbf{Incomplete, net-unidirectional, locus-autonomous sorting.} Robustly,
most diagnostic loci are only partly sorted (\,$\sim$7 of 20 populations fixed) and
unidirectional per locus, with a net aquilonia bias (Module~A); the divergence lives
in \emph{which} populations have sorted \emph{which} loci --- an incidence that is
near-independent across loci (effective dimensionality $\approx15$ of 20; no shared
sorting axis). This locus-autonomy is what a few large selected blocks could not
produce, and it is what cancels the LD.
\item \textbf{Orthogonal to the demographic origin.} Whether the populations are one
connected pool that spread (neutral $F_{ST}\approx0$, our reading) or many
independent hybridizations, the diagnostic-locus pattern is the same --- the neutral
loci decide the origin, the diagnostic loci decide the selection.
\end{itemize}

\section*{Conclusion}
The null is calibrated on the loci where incompatibilities/selection should be absent
(neutral background matches), so the diagnostic-locus behaviour is read against a fair
baseline. There, the data are simultaneously \emph{more differentiated} and \emph{less
linked} than the null --- a combination drift, sampling and panmixia are all forbidden
from producing. The empirical hybrid populations are therefore undergoing
\textbf{divergent, locus-specific, largely intrinsic (DI-governed) sorting of parental
ancestry}, incompletely and heterogeneously across the range. The $F_{ST}$--LD
dissociation of Fig.~\ref{fig:diss} is the most compact evidence for it we have.

\vspace{0.5em}\noindent\footnotesize
Reproduced by \texttt{dev/R/moduleE\_emp\_vs\_sim\_panels.R}; LD is mean $r^2$ to the
directly adjacent SNP(s) throughout; null $=$ 20 independent demes (450/195, $K$=6250,
gen~60) downsampled to empirical population sizes.
\end{document}
